\chapter{Renormalizability and Local Counterterms}
\label{ch:volume-ii-renormalizability-local-counterterms}

\section{Regulated Euclidean Actions}

Let \(\phi\) be a real scalar field on \(D\)-dimensional Euclidean space.
Fix a regulator, denoted collectively by \(\Lambda\).  The regulated
functional integral has the form
\[
  Z_{\Lambda}
  =
  \int [D\phi]_{\Lambda}\,
  \exp\!\left[
    -\int \dd^Dx\,\mathcal L_{E,\Lambda}
  \right],
\]
with bare Euclidean Lagrangian
\[
  \mathcal L_{E,\Lambda}
  =
  \frac12 \partial_{\mu}\phi\,\partial_{\mu}\phi
  +
  \frac12 m^2\phi^2
  +
  \sum_I g_I\,\mathcal O_I(\phi).
\]
Each \(\mathcal O_I\) is a local polynomial, or more generally a local
function, of \(\phi(x)\) and finitely many derivatives
\(\partial_{\mu}\phi(x),\partial_{\mu}\partial_{\nu}\phi(x),\ldots\).
The coefficient \(g_I\) is a bare coupling.  The collection of bare couplings
includes both the parameters chosen as inputs and the counterterm
coefficients needed by the regulator.

Two regulator languages will be used repeatedly.  A momentum cutoff restricts
Fourier modes to
\[
  |\vec k|<\Lambda .
\]
Dimensional regularization computes the same formal integrals in
\[
  D=d-\varepsilon
\]
and removes the regulator by taking \(\varepsilon\to0\).  In both languages
the bare couplings are allowed to depend on the regulator.

The perturbative renormalization problem is the following finite-parameter
problem.  Choose a finite list of physical coordinates
\[
  \lambda_{\rm phys}=(\lambda_1,\ldots,\lambda_N)
\]
on the family of theories under consideration.  These coordinates may be
defined by masses, residues, specified 1PI vertices at a subtraction point, or
other finite observables.  One seeks bare couplings of the form
\[
  g_I=g_I(\lambda_{\rm phys},\Lambda)
  \qquad
  \hbox{or}
  \qquad
  g_I=g_I(\lambda_{\rm phys},\varepsilon)
\]
such that all renormalized quantities being used to define the theory have
finite limits as the regulator is removed.

\section{Renormalized Fields and Effective Couplings}

A convenient sufficient criterion is the finiteness of Euclidean Green
functions of a renormalized field.  Define
\[
  \phi_R(x)=Z_R^{-1/2}\phi(x),
\]
where \(Z_R\) is a regulator-dependent normalization factor.  The criterion is
that, after expressing the bare couplings in terms of
\(\lambda_{\rm phys}\),
\[
  \lim_{\Lambda\to\infty}
  \vev{\phi_R(x_1)\cdots\phi_R(x_n)}_{E,\Lambda}
\]
exists as a distribution for every \(n\) in the class under consideration.
Equivalently, the 1PI effective action expressed as a functional of
\(\phi_R\),
\[
  \widetilde\Gamma_{\Lambda}[\phi_R]
  =
  \Gamma_{\Lambda}[Z_R^{1/2}\phi_R],
\]
has finite local effective couplings.  More explicitly, expand
\[
  \widetilde\Gamma_{\Lambda}[\phi_R]
  =
  S_{\rm kin}[\phi_R]
  +
  \int \dd^Dx\,
  \sum_I g^{\rm eff}_{I,\Lambda}\,
  \mathcal O_I(\phi_R),
\]
where \(S_{\rm kin}\) is the chosen finite quadratic reference action.  The
1PI finiteness criterion is
\[
  \lim_{\Lambda\to\infty}
  g^{\rm eff}_{I,\Lambda}(\lambda_{\rm phys})
  \quad\hbox{finite}
\]
for every local operator coefficient included in the renormalized description.

The same structure appears directly in the bare action.  Suppose the bare
quadratic term is
\[
  S_{\rm quad}[\phi]
  =
  \frac12\int \dd^Dx\,
  \phi(-\partial^2+m_0^2)\phi .
\]
After \(\phi=Z_R^{1/2}\phi_R\), choose a finite mass parameter \(m_R\) and
write
\[
\begin{aligned}
  S_{\rm quad}[Z_R^{1/2}\phi_R]
  &=
  \frac12\int \dd^Dx\,
  \phi_R(-\partial^2+m_R^2)\phi_R     \\
  &\quad
  +
  \frac{Z_R-1}{2}
  \int \dd^Dx\,
  \phi_R(-\partial^2+m_R^2)\phi_R
  +
  \frac12\delta m^2
  \int \dd^Dx\,\phi_R^2 .
\end{aligned}
\]
The first line is the finite quadratic reference action.  The second line is
local and is treated perturbatively as the wavefunction and mass counterterm
part of the interaction.

Thus perturbation theory is organized by writing
\[
  \widetilde S_{\Lambda}[\phi_R]
  =
  S_{\rm ref}[\phi_R]
  +
  S_{\rm int}[\phi_R;\lambda_{\rm phys}]
  +
  S_{\rm ct}[\phi_R;\lambda_{\rm phys},\Lambda],
\]
where \(S_{\rm ct}\) is local and regulator dependent.  The finite parts of
the counterterms are part of the renormalization prescription: changing them
changes the coordinate system \(\lambda_{\rm phys}\) on the same perturbative
family.

\section{Locality as Extension of Distributions}

The locality of counterterms can be formulated without reference to a
particular regulator.  In position space, a Feynman graph with vertices
\[
  x_1,\ldots,x_r
\]
defines a distribution on the configuration space where no ultraviolet
collision has occurred.  Ultraviolet renormalization is the problem of
extending this distribution across the collision diagonals, such as
\[
  x_{i_1}=\cdots=x_{i_k}.
\]
The possible ambiguities in the extension are local distributions supported on
those diagonals.

\begin{definition}[Scaling degree]
Let \(T\in\mathcal D'(\mathbb R^n\setminus\{0\})\).  For \(\lambda>0\), define
the scaled distribution \(T_\lambda\) by
\[
  \langle T_\lambda,f\rangle
  =
  \lambda^{-n}
  \left\langle T, f(\lambda^{-1}\,\cdot)\right\rangle ,
  \qquad f\in C_c^\infty(\mathbb R^n\setminus\{0\}).
\]
The scaling degree of \(T\) at the origin is
\[
  \operatorname{sd}_0(T)
  =
  \inf\left\{\delta\in\mathbb R:
  \lambda^\delta T_\lambda\to0
  \ \hbox{in}\ \mathcal D'(\mathbb R^n\setminus\{0\})
  \ \hbox{as}\ \lambda\to0^+\right\}.
\]
\end{definition}

The extension theorem is the distribution-theoretic core of perturbative
locality.

\begin{theorem}[Extension at a diagonal]
Let \(T\in\mathcal D'(\mathbb R^n\setminus\{0\})\) have finite scaling degree
\(s=\operatorname{sd}_0(T)\).  There exists an extension
\(\overline T\in\mathcal D'(\mathbb R^n)\) whose restriction to
\(\mathbb R^n\setminus\{0\}\) is \(T\) and whose scaling degree is still
\(s\).  If \(s<n\), this extension is unique.  If \(s\ge n\), any two
extensions of scaling degree \(s\) differ by
\[
  \sum_{|\alpha|\le \lfloor s\rfloor-n}
  c_\alpha\,\partial^\alpha\delta^{(n)}(x),
\]
with constants \(c_\alpha\).
\end{theorem}

For a graph subconfiguration collapsing to a diagonal, translation invariance
removes the center-of-mass coordinate, so \(n=D(k-1)\) for \(k\) colliding
vertices.  The integer
\[
  \omega=\lfloor \operatorname{sd}_0(T)\rfloor-n
\]
is the degree of local ambiguity.  A derivative
\(\partial^\alpha\delta^{(n)}(x)\) supported on the diagonal is, after
Fourier transform, a polynomial of degree \(|\alpha|\) in the momenta entering
the collapsed subgraph.  Thus the extension theorem gives exactly the local
Taylor polynomial counterterms used in momentum-space renormalization.

Multiple collision diagonals meet along smaller diagonals.  Choosing
extensions compatibly over this stratified configuration space is the
position-space form of the subdivergence problem.  The forest formulas in the
next chapter organize the same compatibility in momentum space: each divergent
subgraph corresponds to one collision diagonal, and each counterterm is the
allowed derivative-of-delta ambiguity on that diagonal.

\section{First Counterterm Census in Four Dimensions}

The local nature of the required terms is visible in elementary scalar
examples.  In the following diagrams, a shaded disk labelled \(1{\rm PI}\)
denotes the full one-particle-irreducible vertex being tested.

First take a formal \(D=4\) scalar theory with a cubic vertex
\(\int g_3\phi^3\).  The one-loop two-point graph contains the integral
\[
  g_3^2
  \int^{\Lambda}
  \frac{\dd^4q}{(2\pi)^4}\,
  \frac{1}{(q^2+m^2)^2}
  \sim
  g_3^2\log\Lambda .
\]
It contributes to the local operator \(\phi^2\), and is canceled by a mass
counterterm.  The corresponding one-loop triangle and box integrals are
ultraviolet finite in four dimensions.

For \(D=4\) scalar theory with a quartic vertex \(\int g_4\phi^4\), the
two-point 1PI function receives a quadratically divergent tadpole and, at the
next order, a logarithmic momentum-dependent contribution.  The four-point
1PI function has a logarithmic divergence.  These divergences are canceled by
\[
  \delta m^2\phi^2,
  \qquad
  \delta Z\,\partial_{\mu}\phi\,\partial_{\mu}\phi,
  \qquad
  \delta g_4\phi^4 .
\]
Power counting assigns the higher \(n\)-point 1PI functions with \(n\ge6\)
negative superficial degree of divergence, so this interaction closes on a
finite local counterterm list.

For \(D=4\) scalar theory with a sextic vertex \(\int g_6\phi^6\), the first
tests already generate higher local terms.  A one-loop eight-point function
has logarithmic behavior, requiring a local \(\phi^8\) counterterm.  Once
\(\phi^8\) is present, further diagrams generate \(\phi^{10}\), and the
process continues.  A finite number of physical coordinates then does not
determine all local coefficients in the cutoff-removal limit.

\begin{figure}[H]
  \centering
  \resizebox{0.97\textwidth}{!}{%
  \begin{tikzpicture}[>=Stealth,x=1cm,y=1cm,every node/.style={font=\small}]
    \tikzset{
      line/.style={line width=0.95pt},
      blob/.style={circle,draw=black,fill=gray!18,line width=0.9pt,minimum size=0.8cm}
    }
    \begin{scope}
      \node[font=\normalsize] at (0,2.55) {\(D=4,\ \phi^3\)};
      \node[blob] (A) at (-1.4,1.15) {\(1{\rm PI}\)};
      \draw[line] (-2.55,1.15)--(A.west);
      \draw[line] (A.east)--(-0.25,1.15);
      \node at (0.1,1.15) {\(=\)};
      \draw[line] (0.55,1.15)--(2.15,1.15);
      \draw[line] (0.75,1.15)..controls (1.15,1.85) and (1.75,1.85)..(1.95,1.15);
      \node[blue!70!black,align=center] at (1.35,0.3)
        {\(\sim g_3^2\log\Lambda\)\\\(\phi^2\) counterterm};
      \coordinate (u) at (-1.55,-1.1);
      \coordinate (v) at (-0.9,-0.1);
      \coordinate (w) at (-0.25,-1.1);
      \draw[line] (u)--(v)--(w)--(u);
      \draw[line] (v)--(-0.9,0.58);
      \draw[line] (u)--(-2.15,-1.7);
      \draw[line] (w)--(0.35,-1.7);
      \node[green!50!black] at (1.15,-1.1) {triangle finite};
    \end{scope}

    \begin{scope}[shift={(5.6,0)}]
      \node[font=\normalsize] at (0,2.55) {\(D=4,\ \phi^4\)};
      \draw[line] (-2.2,1.15)--(-1.0,1.15);
      \draw[line] (1.0,1.15)--(2.2,1.15);
      \draw[line] (-1.0,1.15)..controls (-0.5,1.8) and (0.5,1.8)..(1.0,1.15);
      \draw[line] (-1.0,1.15)..controls (-0.5,0.5) and (0.5,0.5)..(1.0,1.15);
      \node[blue!70!black,align=center] at (0,0.08)
        {\(\phi^2\) and\\\((\partial\phi)^2\)};
      \draw[line] (-1.45,-1.0)--(1.45,-1.0);
      \draw[line] (-1.45,-1.55)--(1.45,-1.55);
      \draw[line] (-1.45,-1.0)..controls (-1.85,-1.28) and (-1.85,-1.27)..(-1.45,-1.55);
      \draw[line] (1.45,-1.0)..controls (1.85,-1.28) and (1.85,-1.27)..(1.45,-1.55);
      \node[green!50!black] at (0,-2.15) {\(\phi^4\) logarithm};
    \end{scope}

    \begin{scope}[shift={(11.1,0)}]
      \node[font=\normalsize] at (0,2.55) {\(D=4,\ \phi^6\)};
      \node[blob,minimum size=1.0cm] (B) at (0,0.55) {\(1{\rm PI}\)};
      \foreach \a in {20,70,140,200,250,320} {
        \draw[line] ({0.5*cos(\a)},{0.55+0.5*sin(\a)})
          --++({0.82*cos(\a)},{0.82*sin(\a)});
      }
      \draw[line] (-0.32,0.18)..controls (0,-0.45)..(0.32,0.18);
      \node[red!70!black,align=center] at (0,-1.35)
        {higher local terms\\\(\phi^8,\phi^{10},\ldots\)};
    \end{scope}
  \end{tikzpicture}
  }
  \caption{A first counterterm census in four dimensions.  Cubic and quartic
  scalar interactions require local terms from a finite operator list, while a
  sextic interaction generates higher local operators under loop corrections.}
  \label{fig:volume-ii-counterterm-census}
\end{figure}

\section{Engineering Dimensions and the Finite List}

The preceding examples are organized by mass dimension.  The Euclidean action
is dimensionless.  Since the kinetic term is
\[
  \frac12\int \dd^Dx\,\partial_{\mu}\phi\,\partial_{\mu}\phi ,
\]
the field has engineering dimension
\[
  [\phi]=\frac{D-2}{2}.
\]
Let \(\mathcal O_I\) contain \(n_I\) fields and \(\ell_I\) derivatives.  Its
engineering dimension is
\[
  d_I=\ell_I+n_I\frac{D-2}{2},
\]
and the coupling multiplying it has dimension
\[
  [g_I]=D-d_I .
\]

The coefficient \(g_I^{\rm eff}\) of \(\mathcal O_I\) in the 1PI effective
action is extracted by differentiating an \(n_I\)-point 1PI vertex with
respect to the appropriate external momenta and then setting those momenta to
zero, provided the zero-momentum limit is infrared regular.  A diagram built
from inserted vertices \(J\) has the superficial ultraviolet form
\[
  \left(\prod_{J\in G}g_J\right)
  \Lambda^{
    (D-d_I)-\sum_{J\in G}(D-d_J)
  }
  \bigl(\log\Lambda\bigr)^{r},
  \qquad
  0\le r\le L,
\]
where \(L\) is the loop number.  Therefore a counterterm for
\(\mathcal O_I\) can be required only when
\[
  D-d_I-\sum_{J\in G}(D-d_J)\ge0 .
\]

If every coupling appearing in the action has \(d_J\le D\), then
\(D-d_J\ge0\).  The inequality above implies \(d_I\le D\).  For scalar
polynomial interactions with finitely many derivatives at fixed \(D\), there
are only finitely many local operators with \(d_I\le D\).  This is the
power-counting finite-list criterion for perturbative renormalizability.

If an interaction with \(d_J>D\) is included and its coefficient is allowed to
participate in the cutoff-removal problem on the same footing, repeated
insertions can satisfy the inequality for operators of arbitrarily large
dimension.  Such a Lagrangian is still a valid low-energy effective
description at finite cutoff, but it does not define a finite-parameter
cutoff-removal problem by this criterion.

\section{A Renormalizable Model in Six Dimensions}

The scalar cubic interaction is marginal in \(D=6\).  Use the Euclidean
Lagrangian
\[
\begin{aligned}
  \mathcal L_E
  &=
  \frac12\partial_{\mu}\phi\,\partial_{\mu}\phi
  +
  \frac12 m_R^2\phi^2
  +
  \delta Z
  \left(
    \frac12\partial_{\mu}\phi\,\partial_{\mu}\phi
    +
    \frac12 m_R^2\phi^2
  \right)          \\
  &\quad
  +
  \frac12\delta m^2\phi^2
  +
  \frac1{3!}(g_R+\delta g)\phi^3 .
\end{aligned}
\]
The finite parameters are \(m_R\), \(g_R\), and the normalization convention
for \(\phi\).  The local counterterms are
\[
  \delta Z,\qquad \delta m^2,\qquad \delta g .
\]

At order \(g_R^2\), the one-loop self-energy is
\[
  \Sigma_{2}(k)
  =
  \frac12 g_R^2
  \int
  \frac{\dd^D\ell}{(2\pi)^D}\,
  \frac{1}
       {(\ell^2+m_R^2)((k-\ell)^2+m_R^2)} .
\]
Combining denominators gives
\[
  \Sigma_{2}(k)
  =
  \frac12 g_R^2
  \int_0^1\dd x
  \int
  \frac{\dd^D\ell}{(2\pi)^D}\,
  \frac{1}
       {\left(\ell^2+x(1-x)k^2+m_R^2\right)^2}.
\]
In dimensional regularization,
\[
  \int\frac{\dd^D\ell}{(2\pi)^D}\,
  \frac1{(\ell^2+\Delta)^n}
  =
  \frac1{(4\pi)^{D/2}}
  \frac{\Gamma(n-D/2)}{\Gamma(n)}
  \Delta^{D/2-n}.
\]
For \(D=6-\varepsilon\), the pole part of \(\Sigma_2(k)\) is a local
polynomial
\[
  \Sigma_2(k)\big|_{\rm pole}
  =
  A_1(\varepsilon)\,k^2
  +
  B_1(\varepsilon)\,m_R^2,
\]
where \(A_1\) and \(B_1\) have simple \(1/\varepsilon\) poles in minimal
subtraction.  The counterterms \(\delta Z\) and \(\delta m^2\) are chosen so
that the renormalized self-energy is finite.

At order \(g_R^3\), the one-loop three-point triangle has a logarithmic
ultraviolet divergence.  Its pole is independent of external momenta at the
local order relevant for the marginal vertex, so it is canceled by
\(\delta g\).  The one-loop four-point box in \(D=6\) is ultraviolet finite by
power counting.

\begin{figure}[H]
  \centering
  \resizebox{0.95\textwidth}{!}{%
  \begin{tikzpicture}[>=Stealth,x=1cm,y=1cm,every node/.style={font=\small}]
    \tikzset{
      line/.style={line width=0.95pt},
      blob/.style={circle,draw=black,fill=gray!18,line width=0.9pt,minimum size=0.82cm}
    }
    \node[font=\normalsize] at (0,2.0) {\(D=6,\ \phi^3\)};

    \node[blob] (A) at (-5.1,0.75) {\(1{\rm PI}\)};
    \draw[line] (-6.3,0.75)--(A.west);
    \draw[line] (A.east)--(-3.9,0.75);
    \node at (-3.55,0.75) {\(=\)};
    \draw[line] (-3.05,0.75)--(-1.7,0.75);
    \draw[line] (-2.85,0.75)..controls (-2.45,1.38) and (-2.1,1.38)..(-1.9,0.75);
    \node[blue!70!black,align=center] at (-2.35,-0.05)
      {\(A_1 k^2+B_1m_R^2\)\\\(\delta Z,\delta m^2\)};
    \node at (-0.75,0.75) {\(+\)};
    \draw[line] (-0.25,0.75)--(1.05,0.75);
    \draw[red!75!black,line width=1pt] (0.3,0.65)--(0.5,0.85);
    \draw[red!75!black,line width=1pt] (0.3,0.85)--(0.5,0.65);

    \node[blob] (B) at (-4.9,-1.4) {\(1{\rm PI}\)};
    \foreach \a in {90,210,330} {
      \draw[line] ({-4.9+0.41*cos(\a)},{-1.4+0.41*sin(\a)})
        --++({0.85*cos(\a)},{0.85*sin(\a)});
    }
    \node at (-3.55,-1.4) {\(=\)};
    \coordinate (u) at (-2.45,-0.75);
    \coordinate (v) at (-3.0,-1.75);
    \coordinate (w) at (-1.9,-1.75);
    \draw[line] (u)--(v)--(w)--(u);
    \draw[line] (u)--(-2.45,-0.05);
    \draw[line] (v)--(-3.65,-2.25);
    \draw[line] (w)--(-1.25,-2.25);
    \node[green!50!black] at (-0.1,-1.35) {\(\delta g\)};
    \node at (0.65,-1.35) {\(+\)};
    \draw[line] (1.15,-1.4)--(2.15,-1.4);
    \draw[line] (1.65,-1.4)--(1.65,-0.55);
    \draw[red!75!black,line width=1pt] (1.55,-1.5)--(1.75,-1.3);
    \draw[red!75!black,line width=1pt] (1.55,-1.3)--(1.75,-1.5);
  \end{tikzpicture}
  }
  \caption{In six-dimensional \(\phi^3\) theory, the one-loop two-point
  divergence is local in \(k^2\) and \(m_R^2\), and the one-loop three-point
  divergence is local in the cubic vertex.}
  \label{fig:volume-ii-six-dimensional-phi-three-one-loop}
\end{figure}

\section{Locality at Higher Order}

The one-loop renormalized self-energy in the six-dimensional cubic theory has
large Euclidean momentum behavior
\[
  \Sigma_R(q)
  \sim
  (\alpha q^2+\beta)\log\frac{q^2}{\mu^2},
  \qquad
  q^2\gg m_R^2,
\]
for finite constants \(\alpha,\beta\) determined by the subtraction scheme.
When this renormalized subgraph is inserted into a larger self-energy graph,
the ultraviolet region of the remaining integral is represented by
\[
  \int
  \frac{\dd^6q}{(2\pi)^6}\,
  \frac{\Sigma_R(q)}
       {(q^2+m_R^2)^2((k-q)^2+m_R^2)} .
\]
The external momentum dependence is isolated by expanding
\[
\begin{aligned}
  \frac1{(k-q)^2+m_R^2}
  &=
  \frac1{q^2+m_R^2}
  -
  \frac{k^2-2k\cdot q}{(q^2+m_R^2)^2}       \\
  &\quad
  +
  \frac{4(k\cdot q)^2}{(q^2+m_R^2)^3}
  +
  R_3(q,k),
\end{aligned}
\]
where Taylor's formula gives \(R_3(q,k)=O(|k|^3|q|^{-5})\) for \(|q|\gg |k|,m_R\).
Because
\[
  \frac{\Sigma_R(q)}{(q^2+m_R^2)^2}
  =
  O\!\left(\frac{\log q^2}{q^2}\right),
  \qquad |q|\to\infty,
\]
the remainder contributes an integrand
\[
  \dd^6q\;O\!\left(\frac{\log q^2}{|q|^7}\right),
\]
which is ultraviolet convergent.  The displayed Taylor terms are local in the
external momentum through order \(k^2\).  After angular integration, odd powers
of \(q\) vanish and
\((k\cdot q)^2\) becomes proportional to \(k^2q^2\).  Thus the divergent part
contributes only to the \(\phi^2\) and \((\partial\phi)^2\) counterterms.

Dimensional regularization expresses the same fact algebraically.  Integrals
of the form
\[
  \int\frac{\dd^Dq}{(2\pi)^D}
  \frac{\log(q^2+\Delta)}
       {(q^2+\Delta)^n}
  =
  -\frac{\partial}{\partial n}
  \left[
    \frac1{(4\pi)^{D/2}}
    \frac{\Gamma(n-D/2)}{\Gamma(n)}
    \Delta^{D/2-n}
  \right]
\]
produce double and single poles in \(\varepsilon\).  For the self-energy
insertion, the pole part has the local form
\[
  \Sigma_{\rm insertion}(k)\big|_{\rm pole}
  =
  A_{\varepsilon}k^2+B_{\varepsilon},
  \qquad
  A_{\varepsilon}
  =
  \frac{a_2}{\varepsilon^2}+\frac{a_1}{\varepsilon},
  \qquad
  B_{\varepsilon}
  =
  \frac{b_2}{\varepsilon^2}+\frac{b_1}{\varepsilon}.
\]
The coefficients \(a_i,b_i\) are finite constants depending on the finite
renormalization prescription.  The pole part is polynomial in the external
momentum, so it is represented by local counterterms already present in the
six-dimensional cubic Lagrangian.

\begin{figure}[H]
  \centering
  \resizebox{0.98\textwidth}{!}{%
  \begin{tikzpicture}[>=Stealth,x=1cm,y=1cm,every node/.style={font=\small}]
    \tikzset{
      line/.style={line width=0.95pt},
      blob/.style={ellipse,draw=black,fill=gray!18,line width=0.9pt,minimum width=1.1cm,minimum height=0.65cm},
      arr/.style={-{Latex[length=2mm]},line width=0.9pt}
    }
    \begin{scope}
      \node[font=\normalsize,align=center] at (-0.3,2.3)
        {renormalized\\subgraph insertion};
      \draw[line] (-2.9,0.6)--(-1.4,0.6);
      \node[blob] (S) at (-0.8,0.6) {\(\Sigma_R(q)\)};
      \draw[line] (S.east)--(0.5,0.6);
      \draw[line] (0.5,0.6)--(1.5,1.05);
      \draw[line] (0.5,0.6)--(1.5,0.15);
      \node[blue!70!black,align=center] at (-0.55,-0.8)
        {\((\alpha q^2+\beta)\log(q^2/\mu^2)\)};
      \draw[arr] (2.0,0.6)--(3.3,0.6);
    \end{scope}

    \begin{scope}[shift={(5.2,0)}]
      \node[font=\normalsize,align=center] at (0,2.3)
        {Taylor terms\\in \(k\)};
      \draw[line] (-1.65,1.05)--(1.65,1.05);
      \draw[line] (-1.65,0.15)--(1.65,0.15);
      \foreach \x/\lab in {-1.05/{k^0},0/{k^1},1.05/{k^2}} {
        \draw[red!70!black,line width=0.95pt] (\x,-0.05)--(\x,1.25);
        \node[red!70!black] at (\x,-0.35) {\(\lab\)};
      }
      \node[align=center] at (0,-1.1)
        {local pieces determine\\\(\delta m^2\) and \(\delta Z\)};
      \draw[arr] (2.0,0.6)--(3.35,0.6);
    \end{scope}

    \begin{scope}[shift={(10.6,0)}]
      \node[font=\normalsize] at (0,2.3) {finite remainder};
      \draw[line] (-1.4,0.8)..controls (-0.75,1.35) and (0.75,1.35)..(1.4,0.8);
      \draw[line] (-1.4,0.8)..controls (-0.75,0.25) and (0.75,0.25)..(1.4,0.8);
      \draw[line] (-2.4,0.8)--(-1.4,0.8);
      \draw[line] (1.4,0.8)--(2.4,0.8);
      \node[green!45!black,align=center] at (0,-0.55)
        {large-\(q\) falloff\\after subtraction};
    \end{scope}
  \end{tikzpicture}
  }
  \caption{A renormalized subgraph can leave logarithmic large-momentum
  behavior.  Expanding the remaining propagator in the external momentum
  separates the local \(k^0\) and \(k^2\) terms from a convergent remainder.}
  \label{fig:volume-ii-self-energy-insertion-locality}
\end{figure}

\section{Schwinger Parameters and Subdivergent Regions}

The same mechanism appears geometrically in Schwinger parameters.  Consider
the two-loop self-energy graph in \(D=6\) cubic theory with two loop momenta
\(\ell_1,\ell_2\):
\[
  \frac{g_R^4}{2}
  \int
  \frac{\dd^6\ell_1}{(2\pi)^6}
  \frac{\dd^6\ell_2}{(2\pi)^6}
  P(\ell_1)P(k-\ell_1)P(\ell_2)P(k-\ell_2)P(\ell_1-\ell_2),
\]
where
\[
  P(\ell)=\frac1{\ell^2+m_R^2}.
\]
Using
\[
  P(\ell)=\int_0^{\infty}\dd\alpha\,
  \ee^{-\alpha(\ell^2+m_R^2)},
\]
the Gaussian loop integrations produce an integrand of the form
\[
  \frac1{(4\pi)^6}
  \int_0^{\infty}
  \prod_{i=1}^{5}\dd\alpha_i\,
  \frac{
    \ee^{-m_R^2\sum_i\alpha_i-k^2 Q(\alpha)}
  }
  {(\det A)^3},
\]
with
\[
  A=
  \begin{pmatrix}
    \alpha_1+\alpha_2+\alpha_5 & -\alpha_5\\
    -\alpha_5 & \alpha_3+\alpha_4+\alpha_5
  \end{pmatrix}.
\]
The function \(Q(\alpha)\) is nonnegative.  Ultraviolet regions are boundary
regions in which Schwinger parameters of a subgraph go to zero.

There are two one-loop subdivergent regions:
\[
  \alpha_1,\alpha_2,\alpha_5\to0,
  \qquad
  \alpha_3,\alpha_4,\alpha_5\to0 .
\]
They correspond to shrinking the left or right one-loop self-energy subgraph.
The counterterm insertions for the one-loop self-energy subtract the limiting
functions associated with these boundary regions.  After these subtractions,
the remaining overall ultraviolet region is the simultaneous scaling
\[
  \alpha_1,\ldots,\alpha_5\to0 .
\]
Power counting then says that its dependence on the external momentum is at
most quadratic.  Subtracting the \(k^0\) and \(k^2\) Taylor terms leaves a
finite parameter integral.  The overall counterterm is therefore again a
linear combination of \(\phi^2\) and \((\partial\phi)^2\).

\begin{figure}[H]
  \centering
  \resizebox{0.96\textwidth}{!}{%
  \begin{tikzpicture}[>=Stealth,x=1cm,y=1cm,every node/.style={font=\small}]
    \tikzset{line/.style={line width=0.95pt}}
    \begin{scope}
      \node[font=\normalsize] at (0,2.15) {two-loop self-energy};
      \draw[line] (-3.0,0.5)--(-2.1,0.5);
      \draw[line] (-2.1,0.5)..controls (-1.35,1.35) and (-0.45,1.35)..(0.0,0.5);
      \draw[line] (-2.1,0.5)..controls (-1.35,-0.35) and (-0.45,-0.35)..(0.0,0.5);
      \draw[line] (0.0,0.5)..controls (0.45,1.35) and (1.35,1.35)..(2.1,0.5);
      \draw[line] (0.0,0.5)..controls (0.45,-0.35) and (1.35,-0.35)..(2.1,0.5);
      \draw[line] (2.1,0.5)--(3.0,0.5);
      \draw[orange!85!black,dashed,line width=1pt] (-2.3,-0.52) rectangle (0.18,1.52);
      \draw[green!55!black,dashed,line width=1pt] (-0.18,-0.52) rectangle (2.3,1.52);
      \node[orange!85!black] at (-1.18,-0.9)
        {\(\alpha_1,\alpha_2,\alpha_5\to0\)};
      \node[green!55!black] at (1.18,-0.9)
        {\(\alpha_3,\alpha_4,\alpha_5\to0\)};
    \end{scope}

    \begin{scope}[shift={(7.7,0)}]
      \node[font=\normalsize] at (0,2.15) {nested subtraction};
      \draw[line] (-2.1,0.7)..controls (-1.2,1.45) and (1.2,1.45)..(2.1,0.7);
      \draw[line] (-2.1,0.7)..controls (-1.2,-0.05) and (1.2,-0.05)..(2.1,0.7);
      \draw[line] (-3.0,0.7)--(-2.1,0.7);
      \draw[line] (2.1,0.7)--(3.0,0.7);
      \foreach \x in {-1.3,1.3} {
        \draw[red!70!black,line width=1pt] (\x,-0.25)--(\x,1.45);
      }
      \node[red!70!black,align=center] at (0,-0.75)
        {subgraphs first};
      \node[purple!80!black,align=center] at (0,-1.45)
        {overall \(k^0,k^2\) terms last};
    \end{scope}
  \end{tikzpicture}
  }
  \caption{In Schwinger parameters, ultraviolet divergences are boundary
  regions.  Subdivergent one-loop regions are subtracted before the remaining
  overall boundary region is expanded in local external-momentum terms.}
  \label{fig:volume-ii-schwinger-subdivergences}
\end{figure}

The preceding analysis is the local core of perturbative renormalizability for
power-counting finite-list theories.  At all orders, the same nested structure
is organized by the Bogoliubov--Parasiuk--Hepp--Zimmermann theorem: subtract
local Taylor polynomials for divergent subgraphs in a compatible nested order,
and the remaining integrals are finite.  The next chapter develops this
subtraction structure as a systematic forest formula.
